\chapter{オプショナル値}

wick に存在理由が一つだけあるとすれば、オプショナル値です。静的な型、真偽の暗黙変換の禁止、コンパイル時のエンジン呼び出し検査は、すべて一つの主張を支えるものです。「この値は存在しないかもしれない」という事実は、原因から三つ先の関数で突然現れる実行時の問題ではなく、コンパイラが知り、プログラマが処理すべき事実である。本章ではオプショナル値を詳しく説明し、それを譲れない要件にしたバグの話を紹介します。

\section{Lua で起きた問題}

失敗から始めましょう。この機能は、その失敗の形にぴったり合わせて設計したからです。Lantern Night は保存ファイルに最高得点を記録します。Lua 版は次のように読み込んでいました。

\begin{lstlisting}[language=luatwin]
local best = tonumber(lt.load_save("lanternnight_best") or "0")
\end{lstlisting}

\noindent
どの部分も妥当に見えます。`lt.load_save` は保存した文字列を返し、保存データがなければ `nil` を返します。`or "0"` は `nil` のときの既定値を与え、`tonumber` は文字列を数値として解析します。正しく読め、レビューを通り、出荷されました。

その後、プレイヤーが保存中にゲームを強制終了しました。保存はディスクへの書き込みです。まずい瞬間に中断すると、ファイルだけが作られ、中身が空の0バイトになります。次の起動時、`lt.load_save` はそのファイルを開き、正常に読み取り、`nil` ではなく空文字列 `""` を返しました。`""` を手に、式をもう一度たどってみましょう。

\begin{itemize}
\item `lt.load_save(...)` は `""` です。`nil` ではなく、実在する文字列です。
\item Lua では空文字列が\emph{真として扱われる}ので、`"" or "0"` は `""` です。代替値は使われません。
\item 空文字列は数値ではないため、`tonumber("")` は `nil` です。
\item `local best = nil` となり、以後の `best` の比較、算術、数値の描画は、すべて `nil` に対して行われます。
\end{itemize}

\noindent
ゲームは最初のフレームでクラッシュしました。保存や読み込みの場所ではなく、その先で、`nil` が初めて許容されない演算に出会ったところです。スタックトレースは、何の落ち度もない行を指していました。この現象には Lua の四つの判断が重なる必要がありました。空文字列を真とすること、空ファイルの読み込みが `nil` ではなく `""` を返すこと、解析失敗を `nil` で表すこと、`nil` の伝播を止める静的検査がないことです。それぞれ単独では合理的です。それでも組み合わさると、ディスク書き込みの中断が、コードを読んでも見抜けない最初のフレームのクラッシュに変わりました。

\section{バグを書けなくする型}

wick の答えは型です。値を生成できない可能性のある操作は、その値を直接返すのではなく、`T?` と書く\emph{オプショナル値}を返します。これは `T` または `nil` です。コンパイラは `T?` を `T` として使うことを許しません。

\begin{lstlisting}[language=wick]
let s = lt.load_save("hi")       // s: str?
lt.print(s, 4, 4, 1, 1, 1, 1)    // COMPILE ERROR: expected str, got str?
\end{lstlisting}

\noindent
二行目はコンパイルできません。エラーは \texttt{expected str, got str?} で、その行を書いた瞬間に、プレイヤーの端末ではなく自分の端末で現れます。`str` を必要とする `lt.print` に `str?` が届く経路はありません。`str?` から `str` に進むには、値が `nil` のときどうするかをプログラム内で示す必要があります。wick にはその方法が二つあります。

\section{\texttt{??} で中の値を取り出す}

nil 合体演算子 `??` は、`nil` の場合に既定値を与えて、オプショナル値から中の値を取り出します。`a ?? b` は `a` が `nil` でなければ `a`、そうでなければ `b` になります。結果の型は非オプショナルの元の型です。

\begin{lstlisting}[language=wick]
let text = s ?? "no save yet"    // text: str
let n = num("maybe") ?? 0        // n: num  (num parses str -> num?)
\end{lstlisting}

\noindent
既定値は元の型に一致しなければなりません。`s` は `str?`、既定値は `str` なので `s ?? "..."` は使えます。一致しないと \texttt{'??' default must be str} と報告されます。str の部分は元の型に応じて変わります。右辺は左辺が `nil` のときだけ遅延評価されるので、そこで処理を行っても安全です。また、`??` の左辺は実際にオプショナル値でなければなりません。`nil` になり得ない値に使うと、合体させる必要がないため \texttt{'??' left side must be an optional} というエラーになります。

実例に戻りましょう。あの致命的な Lua の行を wick にすると、次のようになります。

\begin{lstlisting}[language=wick]
let best = num(lt.load_save("lanternnight_wick_best") ?? "") ?? 0
\end{lstlisting}

\noindent
Lua 版をクラッシュさせた0バイトのファイルをたどります。

`lt.load_save(...)` は `str?` を返します。ファイルがなくて `nil` なのか、空ファイルで `""` なのかは、その後の処理に影響しません。最初の `??` が\emph{どちらも} `str` にするからです。`nil` は `""` になり、`""` はすでに `""` です。次の `num("")` は `num?` で、空文字列は数値ではないので `nil` です。二つ目の `??` が受け止め、`0` を返します。保存ファイルが存在するか、形式が正しいか、中身が空かによらず、結果は必ず普通の `num` です。バグを直したというより、バグを\emph{書けなくした}のです。どちらかの `??` を消すと、コンパイルできません。

この違いを身につけることが大切です。Lua では安全な版とバグのある版がほぼ同じに見え、コンパイラは両方を受け付けます。違いが現れるのは、本番環境で特定のまれな入力が来たときだけです。wick では `??` が欠けたバグのある版はコンパイルエラーで、安全な版だけがビルドできます。型システムは、バグの発見をプレイヤーの最初のフレームから、作者の最初の保存へ移したのです。

\section{\texttt{if x != nil} で型を絞り込む}

オプショナル値を扱う二つ目の方法は、型を\emph{絞り込む}ことです。`nil` と比較し、検査を通過したブロック内で普通の `T` として使います。

\begin{lstlisting}[language=wick]
if s != nil {
  // inside this block, s has type str, not str?
  lt.print(s, 4, 4, 1, 1, 1, 1)
}
\end{lstlisting}

\noindent
`s != nil` は二つの役割を果たします。`bool` の条件であると同時に、保護されたブロック内では `s` が `nil` でないことを型検査器に伝えるので、その中の型は `str` になります。読む人が直感的に行う「nil でないと確かめたので、ここでは文字列だ」という判断を、形式化し強制しています。

値を使っていくつかの処理を行いたく、既定値では扱いにくいときや、`nil` の場合に独自の処理があるときは、絞り込みが適しています。

\begin{lstlisting}[language=wick]
let saved = lt.load_save("name")
if saved != nil {
  lt.print("WELCOME BACK " + saved, 4, 4, 1, 1, 1, 1)
} else {
  lt.print("NEW PLAYER", 4, 4, 1, 1, 1, 1)
}
\end{lstlisting}

\paragraph{絞り込みの制約。} v0.3 の絞り込みは\emph{ローカル変数}に適用されます。モジュールのグローバル変数は、検査と使用の間に別の関数や後のフレームのコードが変更し得るため、`if` では絞り込まれません。グローバルのオプショナル値には `??` を使うか、先にローカル変数へコピーしてから絞り込みます。ブロック内で、絞り込んだ変数に新しいオプショナル値を代入すると、絞り込みは\emph{解除}されます。保証があるのは検査した値だけです。また、v0.3 では `and`・`or` の連なりを通じて絞り込みが伝わりません。組み合わせた条件にコンパイラが括弧を求める理由です（第3章）。迷ったら入れ子の `if` で絞り込むか、ローカル変数へコピーします。いずれも第8章に記載した既知の v0.3 の制約で、バグではありません。

\section{オプショナル値を返す操作}

オプショナル値は適当に付け足すものではありません。実際に失敗し得る操作から生まれ、コンパイラがそこから伝播させます。v0.3 では以下が該当します。

\begin{center}
\begin{tabular}{@{}llp{0.40\linewidth}@{}}
\toprule
\textbf{式} & \textbf{型} & \textbf{\texttt{nil} になる場合} \\
\midrule
\texttt{lt.load\_save(k)} & \texttt{str?} & 読み込める保存データがない \\
\texttt{num(s)} & \texttt{num?} & \texttt{s} が数値でない（\texttt{""} も含む） \\
\texttt{pop(xs)} & \texttt{T?} & リストが空 \\
\texttt{m[key]} & \texttt{T?} & キーが存在しない \\
\texttt{env(name)} & \texttt{str?} & 変数が設定されていない \\
\bottomrule
\end{tabular}
\end{center}

\noindent
いずれも他の言語なら、空文字列、偽、`nil`、「見つからない」を表す特殊な値など、確認を忘れやすい番兵値を受け取るところです。wick は失敗を型で表すため、確認を忘れることができません。忘れた確認はコンパイルエラーになります。最も分かりやすいのはマップの検索です。`m[key]` は、キーの存在を確信していても\emph{必ず} `T?` です。型システムには、あなたの確信は分からないからです。キーがない場合を呼び出し箇所で必ず意識することになり、`??` なら三文字で処理できます。

\begin{lstlisting}[language=wick]
let hp = stats["health"] ?? 100
\end{lstlisting}

\section{型システムの中のオプショナル値}

オプショナル値と言語の他の部分をつなぐ規則を、第3章から、オプショナル値の側から見直してまとめます。

\begin{itemize}
\item \textbf{拡張はコストがありません。} `T?` が必要な場所には `T` を使えます。「必ずある」から「ないかもしれない」へ移っても、安全性は失われません。
\item \textbf{中の値を取り出す必要があります。} `T` が必要な場所に `T?` は使えません。移れるのは `??` と絞り込みだけで、どちらも `nil` を扱う必要があります。
\item \textbf{オプショナル値は入れ子になりません。} `T??` はありません。オプショナル値は `T` または `nil`、それだけです。
\item \textbf{\texttt{nil} には型が必要です。} 単独の `nil` からは型を推論できません。`let x: str? = nil` と書けば指定できます。`nil` は、特定のオプショナル型の空の場合としてだけ意味を持ちます。
\item \textbf{非オプショナル値と \texttt{nil} の比較はエラーです。} `x` が普通の `str` なら、`x` は `nil` になり得ないので `x != nil` は無意味なコードです。コンパイラは必ず成功する検査を許すのではなく、\texttt{comparing non-optional to nil} と伝えます。
\end{itemize}

\section{実際の書き方}

日々 wick を使うと、動的な言語が要求する防御的な足場を書かなくなります。足場が守っていた事実を、コンパイラが追跡するからです。読み込みと解析を毎回 `if x ~= nil` の連なりで囲む必要はありません。値の入り口で一度 `??` を書けば、その先は裏切らない普通の `T` です。オプショナル値は余計な負担ではありません。どの値が `nil` かもしれないか覚えておく、手作業で、忘れやすく、レビューしても見逃す負担を取り除きます。代わりにコンパイラが覚え、忘れることを拒みます。

0バイトの保存ファイルは一つのバグです。しかし「失敗し得る操作の失敗を普通の値で表せるため、確認を省けてしまう」という種類のバグは膨大です。wick はこの種類をなくすために作られました。wick プログラムにある一つ一つの `T?` は、その種類の一例を処理したという小さな証明です。
